\chapter{集合、函数与证明}
\label{chap:sets-functions-and-proofs}

机器学习中的公式往往写得很短。例如，给定样本
$(\bm{x}_i,y_i)\in\mathbb{R}^d\times\mathbb{R}$，其中$i=1,\ldots,n$，并取
正则化系数$\lambda\geq0$，一个带正则项的最小二乘问题可以写成
\begin{equation}
  \min_{\bm{w}\in\mathbb{R}^d}
  F(\bm{w})
  =
  \frac{1}{2n}\sum_{i=1}^{n}
  \bigl(\bm{w}^{\top}\bm{x}_i-y_i\bigr)^2
  +
  \frac{\lambda}{2}\lVert\bm{w}\rVert_2^2.
  \label{eq:prelim-regularized-least-squares}
\end{equation}
其中$\lambda=0$对应不加正则项的情形。这行公式同时规定了候选参数、样本索引、目标函数和比较准则。若省略
$\bm{w}\in\mathbb{R}^d$，读者便不知道优化在哪个集合上进行；若不说明
$\lambda$的范围，正则项可能鼓励参数变大而不是抑制参数；若把“存在一个参数对所有
样本都预测正确”误写成“对每个样本都存在一个预测正确的参数”，结论甚至会完全改变。

本章把这些容易隐没在公式里的约定展开。集合说明对象从哪里取值，函数说明对象之间怎样
对应，索引把一族对象组织起来，量词规定结论覆盖多大范围，证明则解释结论为何由条件推出。
在这些语言之上，我们才能准确区分最优值、最优解、近似误差和计算代价。后续章节会增加
分析、代数与概率工具，但它们都依赖本章建立的表达方式。

\section{集合、关系与函数}
\label{sec:prelim-sets-relations-functions}

\subsection{从元素到结构}

\term{集合}（set）把若干对象作为一个整体讨论。若对象$x$属于集合$X$，记作
$x\in X$；若$X$的每个元素也属于$Y$，记作$X\subseteq Y$。空集记为
$\varnothing$。两个集合的并、交与差分别写作
\[
  X\cup Y,\qquad X\cap Y,\qquad X\setminus Y.
\]
这些运算不仅整理对象，还能表达条件。例如，若$C$是满足资源约束的参数集合，
$D$是使模型有定义的参数集合，那么真正的可行域是$C\cap D$，而不是其中任意一个集合。

集合的描述常有两种方式。列举法直接写出元素，如
$A=\{1,2,3\}$；条件法写成
\[
  A=\{i\in\mathbb{N}:1\leq i\leq 3\}.
\]
冒号右侧是成员必须满足的条件。集合只记录元素是否属于其中，不记录排列顺序，也不重复
计数。因此$\{1,2,3\}$、$\{3,2,1\}$和$\{1,1,2,3\}$表示同一个集合。需要顺序或重复
次数时，应使用序列、多重集合或向量，而不能继续依赖普通集合。

两个集合$X$与$Y$的\term{笛卡尔积}（Cartesian product）是所有有序对组成的集合
\[
  X\times Y=\{(x,y):x\in X,\ y\in Y\}.
\]
“有序”意味着$(x,y)$与$(y,x)$通常不是同一个对象。一个带标签样本可以表示为
$(\bm{x},y)\in\mathbb{R}^d\times\mathbb{R}$；包含$n$个样本的数据集则可以先看成
该集合中的一个长度为$n$的序列。这样写清了一个样本的组成，却还没有对样本怎样产生
作任何概率假设。

集合$X$上的\term{幂集}（power set）$\mathcal{P}(X)$由$X$的全部子集组成。
如果$X=\{a,b\}$，那么
\[
  \mathcal{P}(X)=\{\varnothing,\{a\},\{b\},\{a,b\}\}.
\]
幂集在后文有两类用途。在概率模型中，事件是样本空间$\Omega$的子集；实际允许赋予概率
的事件通常只组成$\mathcal{P}(\Omega)$的一个结构化子族，有限或可数离散样本空间中常取
完整幂集。幂集也可以描述算法可能选择的特征子集或边集合。并非每次出现“若干子集”都
需要完整幂集；允许的子集受到结构限制时，应另写一个集合族。概率事件族需要满足的闭包
条件将在概率论部分说明。

\subsection{关系把对象联系起来}

集合$X$与$Y$之间的\term{二元关系}（binary relation）是
$X\times Y$的一个子集$R$。写$xRy$表示$(x,y)\in R$。关系允许一个$x$对应多个
$y$，也允许没有对应对象。例如，“不大于”是$\mathbb{R}$上的关系，“样本$i$与样本$j$
之间有边”是索引集合上的关系。

定义在同一个集合$X$上的关系若满足
\begin{itemize}
  \item 对每个$x\in X$都有$xRx$；
  \item $xRy$蕴含$yRx$；
  \item $xRy$且$yRz$蕴含$xRz$，
\end{itemize}
就称为\term{等价关系}（equivalence relation）。三条性质分别称为自反性、对称性与
传递性。等价关系把$X$分成互不相交的\term{等价类}（equivalence class）。例如，在
$\mathbb{R}^2\setminus\{\bm{0}\}$上规定$\bm{x}\sim\bm{y}$当且仅当存在
$c\in\mathbb{R}\setminus\{0\}$使$\bm{x}=c\bm{y}$，那么同一条过原点直线上的非零
向量属于同一个等价类。
后文讨论参数不可识别时，不同参数也可能因为产生同一个可观测对象而落入同一等价类。

\term{偏序关系}（partial order）也要求自反性与传递性，但把对称性换成反对称性：
$xRy$且$yRx$时必须有$x=y$。集合包含关系$\subseteq$就是偏序。偏序中的两个元素
未必可比较，例如$\{1\}$和$\{2\}$都不是对方的子集。这一点会在多目标优化中再次出现：
两个方案可能各自在一个目标上更好，因而不能排成一条总的优劣顺序。

\subsection{函数、复合与逆}

\term{函数}（function）$f:X\to Y$为每个$x\in X$指定唯一的$f(x)\in Y$。
$X$称为定义域，$Y$称为陪域。函数在定义域上的实际输出组成
\[
  f(X)=\{f(x):x\in X\},
\]
称为像。陪域是函数定义的一部分，像则由函数实际取到的值决定，二者不必相等。
例如$f:\mathbb{R}\to\mathbb{R}$、$f(x)=x^2$的像是$[0,\infty)$。
给定$B\subseteq Y$，集合
\[
  f^{-1}(B)=\{x\in X:f(x)\in B\}
\]
称为$B$在$f$下的\term{原像}（preimage）。这里的$f^{-1}(B)$只表示满足输出条件的
输入集合，即使$f$没有逆函数也有定义；它不能与后文双射的逆函数混为一谈。

必须区分函数$f$、输入$x$与函数值$f(x)$。式
\eqref{eq:prelim-regularized-least-squares}中的$F$是把每个参数向量映射到一个实数的
函数，$F(\bm{w})$才是参数$\bm{w}$对应的目标值。若要讨论一组候选函数，应把它们写成
函数集合，例如
\[
  f_{\bm{w}}:\mathbb{R}^d\to\mathbb{R},
  \qquad
  f_{\bm{w}}(\bm{x})=\bm{w}^{\top}\bm{x},
  \qquad
  \mathcal{F}=\{f_{\bm{w}}:\bm{w}\in\mathbb{R}^d\}.
\]
这里$\bm{w}$索引函数，$\bm{x}$是函数的输入。混淆这两个角色，会使“优化参数”和
“对新输入求值”看起来像同一步操作。
更准确地说，参数化给出了映射$\bm{w}\mapsto f_{\bm{w}}$。它的输入是参数向量，输出
却是一个函数，而不是这个函数在某个固定输入上的函数值；这里的函数值是标量。这个映射
若不是单射，不同参数就会表示同一个函数；因此参数集合与由它表示的函数集合也不能直接
视为同一个集合。

函数$f:X\to Y$若满足$f(x_1)=f(x_2)$必然推出$x_1=x_2$，就称为
\term{单射}（injective function）；它不会把两个不同输入压成同一个输出。若每个
$y\in Y$都能写成$f(x)$，就称为\term{满射}（surjective function）；它覆盖整个陪域。
同时为单射和满射的函数称为双射。只有双射才存在定义在整个$Y$上的逆函数
$f^{-1}:Y\to X$，并满足
\[
  f^{-1}(f(x))=x,\qquad f(f^{-1}(y))=y.
\]
若$f$只是单射，可以在像$f(X)$上定义逆；若$f$不是单射，同一个输出对应多个输入，
所谓“反求参数”便需要额外约束或只能返回一个集合。

给定$f:X\to Y$与$g:Y\to Z$，它们的\term{复合}（composition）是
\[
  (g\circ f)(x)=g(f(x)).
\]
复合次序不能任意交换，因为$g\circ f$与$f\circ g$的定义域可能不同，即使都能定义，
结果也通常不同。机器学习模型常由多层变换复合而成；检查每层输入输出的集合或维度，
是发现公式错误的一种直接方法。
再给定$h:Z\to U$，复合满足结合律：
\[
  h\circ(g\circ f)=(h\circ g)\circ f.
\]
恒等映射$\operatorname{id}_X(x)=x$不改变输入，并满足
$f\circ\operatorname{id}_X=f$与$\operatorname{id}_Y\circ f=f$。若$f$与$g$均为
双射，则$g\circ f$也是双射，且$(g\circ f)^{-1}=f^{-1}\circ g^{-1}$；逆映射的次序
与原复合次序相反。

\begin{example}[参数不同而函数相同]
\label{ex:prelim-nonidentifiable-parameterization}
考虑参数化函数
\[
  f_{a,b}:\mathbb{R}\to\mathbb{R},
  \qquad
  f_{a,b}(x)=abx,
  \qquad
  (a,b)\in\mathbb{R}^2.
\]
参数$(1,2)$与$(2,1)$不同，却对每个$x$给出相同的函数值。更一般地，只要$c\neq0$，
$(ca,b/c)$与$(a,b)$表示同一个函数。由于定义域包含非零输入，例如取$x=1$，从完整的
输入输出关系可以确定乘积$ab$，却不能分别确定$a$和$b$。这是参数化的不可识别性，而
不是计算精度不足。
\end{example}

\section{索引、有限和与递推}
\label{sec:prelim-indices-sums-recurrences}

\subsection{同一个下标可能承担不同角色}

索引把一族对象排成可引用的形式。本书通常用$i,j,k$标识元素或样本，用$n$表示样本数，
用$d$表示特征维数，用带括号的上标$t$表示迭代步。例如，
$\bm{x}_i\in\mathbb{R}^d$是第$i$个样本，
$x_{i,j}$是该样本的第$j$个坐标，
$\bm{w}^{(t)}$是第$t$次迭代的参数。普通上标则保留给幂：
$w_j^2$表示平方，而$w_j^{(2)}$表示第$2$次迭代中的第$j$个坐标。
若对象本身按时间排列，还应另行声明时间下标，例如$\bm{z}_{\tau}$表示时刻$\tau$的
状态。样本编号、物理时间和算法迭代都可以取整数值，但它们索引的是不同对象，不能仅因
数值相同就互换。

下标只有在取值范围明确时才有意义。写
\[
  \sum_{i=1}^{n}a_i
\]
表示把$a_1,\ldots,a_n$相加，其中$i$是求和过程中变化的\term{哑索引}
（dummy index）。把$i$换成$k$不会改变结果。与之不同，等式
\[
  b_j=\sum_{i=1}^{n}x_{i,j}
\]
中的$j$是自由索引：等式对哪些$j$成立需要在附近说明，而且改变$j$会选取不同坐标。
同一个符号不能在一项中既作哑索引又作自由索引。

双重求和需要同时写清两个范围。对有限数组$a_{i,j}$，
\[
  \sum_{i=1}^{n}\sum_{j=1}^{d}a_{i,j}
  =
  \sum_{j=1}^{d}\sum_{i=1}^{n}a_{i,j},
\]
因为两边都遍历同一个有限索引集
$\{1,\ldots,n\}\times\{1,\ldots,d\}$，并把每个索引对$(i,j)$对应的项$a_{i,j}$
恰好相加一次。后文把有限和换成无限级数或积分时，交换次序将需要额外条件，不能直接沿用
这个有限情形的结论。

\subsection{有限和、乘积与加权平均}

有限和满足线性关系。对标量$a,b$与同长度序列$(u_i)$、$(v_i)$，
\[
  \sum_{i=1}^{n}(au_i+bv_i)
  =
  a\sum_{i=1}^{n}u_i+b\sum_{i=1}^{n}v_i.
\]
若权重$\alpha_i\geq0$且$\sum_{i=1}^{n}\alpha_i=1$，则
$\sum_i\alpha_i u_i$称为有限\term{加权平均}（weighted average）。
普通样本平均对应$\alpha_i=1/n$。这个简单对象会成为积分与期望的有限版本：积分允许
位置形成连续集合，期望则按概率分配权重。

有限乘积写作$\prod_{i=1}^{n}a_i$。按照空操作不改变结果的约定，
\[
  \sum_{i\in\varnothing}a_i=0,\qquad
  \prod_{i\in\varnothing}a_i=1.
\]
这两个约定分别把$0$和$1$作为加法与乘法的单位元，使递推公式在边界处仍可保持统一。
使用乘积前仍应检查零因子与符号；对正数取对数可把乘积变成和，但含零或负数时不能直接
这样做。

\subsection{递推式描述状态怎样更新}

\term{递推}（recurrence）用当前或过去的值定义后续值。仅写更新规则还不够，必须同时
给出初值。例如
\begin{equation}
  e^{(t+1)}=\rho e^{(t)},\qquad e^{(0)}=e_0,
  \label{eq:prelim-geometric-recurrence}
\end{equation}
逐步展开得到
\[
  e^{(t)}=\rho^t e_0.
\]
当$|\rho|<1$时，误差绝对值按几何速度缩小；$\rho=-1$时大小不变而符号交替；
$|\rho|>1$时则放大。因而“反复执行更新”本身不保证收敛，更新系数的范围才决定长期行为。

更一般的误差递推
\begin{equation}
  e^{(t+1)}\leq \rho e^{(t)}+\delta_t,
  \qquad 0\leq\rho<1,
  \label{eq:prelim-perturbed-recurrence}
\end{equation}
包含每一步的扰动$\delta_t$。连续展开$t$步可得
\begin{equation}
  e^{(t)}
  \leq
  \rho^t e^{(0)}
  +
  \sum_{s=0}^{t-1}\rho^{t-1-s}\delta_s.
  \label{eq:prelim-unrolled-recurrence}
\end{equation}
第一项是初始误差的衰减，第二项是历次扰动经过不同程度衰减后的累计。只说明
$\delta_t$很小，仍不足以断言累计项趋于零；还要检查扰动是否趋零、是否可求和，以及
几何权重怎样作用于它。

式~\eqref{eq:prelim-unrolled-recurrence}可以用数学归纳法证明：先核对$t=1$的基例，再
假设结论对$t$成立并据此推出$t+1$时的结论。完整证明将在本章介绍归纳法时给出。

\subsection{计数公式说明候选有多少}

从$n$个不同对象中依次取出$k$个且不放回，若顺序重要，可得到
\[
  P(n,k)=n(n-1)\cdots(n-k+1)=\frac{n!}{(n-k)!}
\]
种排列。若只关心选中了哪些对象，不关心次序，每个大小为$k$的子集被计算了$k!$次，
因此组合数为
\begin{equation}
  \binom{n}{k}=\frac{n!}{k!(n-k)!},
  \qquad 0\leq k\leq n.
  \label{eq:prelim-binomial-coefficient}
\end{equation}
规定$\binom{n}{k}=0$用于$k<0$或$k>n$，可以统一许多边界公式。

把$n$个对象的子集按是否包含某个固定对象分组，可得Pascal恒等式
\[
  \binom{n}{k}
  =
  \binom{n-1}{k}
  +
  \binom{n-1}{k-1}.
\]
第一项计算不包含固定对象的子集，第二项先包含该对象，再从其余$n-1$个对象中选择
$k-1$个。这种“为同一集合建立不重不漏的两种计数”称为组合证明。后文的组合数学章会系统
讨论计数、生成函数与离散结构；这里保留最常用的公式与论证方式。

二项式定理把组合数与代数展开连接起来：
\begin{equation}
  (a+b)^n
  =
  \sum_{k=0}^{n}\binom{n}{k}a^k b^{n-k}.
  \label{eq:prelim-binomial-theorem}
\end{equation}
选择展开式中的一项时，需要从$n$个因子中挑出$k$个贡献$a$，其余贡献$b$，所以
$a^k b^{n-k}$前的系数正是$\binom{n}{k}$。概率论中的二项分布和学习理论中的有限类
计数都会复用这个关系。

\section{命题、量词与条件}
\label{sec:prelim-propositions-quantifiers}

\subsection{命题必须能够判断真假}

\term{命题}（proposition）是能够判定真假的陈述。“$F$有最小值”是命题，
“求$F$的最小值”则是任务。带变量的表达式如“$x^2\geq0$”还需要说明$x$的取值范围
或添加量词，才成为完整命题。对实数它成立，对复数若使用通常的序关系则没有定义。

命题$P$的否定记为$\neg P$。两个命题的“且”“或”分别记为
$P\land Q$和$P\lor Q$。数学中的“或”通常是可兼容的：$P$和$Q$同时为真时，
$P\lor Q$仍为真。若只允许恰有一个成立，需要明确写成异或条件。

\term{蕴含}（implication）$P\Rightarrow Q$表示只要$P$为真，$Q$就必须为真。
这里$P$是$Q$的充分条件，$Q$是$P$的必要条件。反向命题$Q\Rightarrow P$需要单独
证明。两向都成立时写作$P\Leftrightarrow Q$，称为等价条件。

定理中的假设正是在限定$P$所覆盖的对象。例如，“目标函数连续且可行域紧致”只允许把
结论用于同时满足这两个条件的问题，不能从定理名称中删去其中一个条件。类似地，“存在
一个最优解”是$\exists x^\star\in X$开头的命题，而“每个初值产生的迭代都收敛”是
$\forall x^{(0)}\in X$开头的另一命题；前者成立并不自动给出后者。

以实数$x$为例，$x>1$是$x^2>1$的充分条件，却不是必要条件，因为$x<-1$时也有
$x^2>1$。条件$x\neq0$是$x^2>0$的充分必要条件。把充分条件误当作必要条件，会排除
本来可能成立的情形；把必要条件误当作充分条件，则可能把不合格对象当成解。

\subsection{全称与存在量词规定保证范围}

\term{全称量词}（universal quantifier）$\forall$表示“对每一个”，
\term{存在量词}（existential quantifier）$\exists$表示“至少存在一个”。例如
\[
  \forall x\in\mathbb{R},\quad x^2\geq0
\]
是全称命题，而
\[
  \exists x\in\mathbb{R},\quad x^2=2
\]
是存在命题。存在并不表示唯一；要表达恰有一个对象满足条件，可以写
$\exists!x$，但使用时仍应说明存在性与唯一性各由什么保证。

否定量词时，量词种类与内部命题同时改变：
\begin{align*}
  \neg\bigl(\forall x\in X,\ P(x)\bigr)
  &\Longleftrightarrow
  \exists x\in X,\ \neg P(x),\\
  \neg\bigl(\exists x\in X,\ P(x)\bigr)
  &\Longleftrightarrow
  \forall x\in X,\ \neg P(x).
\end{align*}
所以，要推翻“每个输入都满足性质$P$”，找到一个不满足的输入即可；要证明“没有输入
满足$P$”，则必须排除所有输入。

\subsection{量词次序不能交换}

考虑样本集合$S=\{1,\ldots,n\}$和参数集合$W$，令$C(w,i)$表示参数$w$在第$i$个
样本上预测正确。陈述
\begin{equation}
  \exists w\in W,\quad \forall i\in S,\quad C(w,i)
  \label{eq:prelim-one-parameter-all-samples}
\end{equation}
要求找到同一个参数，使它对所有样本都正确。陈述
\begin{equation}
  \forall i\in S,\quad \exists w\in W,\quad C(w,i)
  \label{eq:prelim-each-sample-some-parameter}
\end{equation}
却允许参数随样本改变，因而弱得多。

\begin{example}[逐点可选不等于统一可选]
\label{ex:prelim-quantifier-order}
令$S=W=\{0,1\}$，并规定$C(w,i)$当且仅当$w=i$。对每个$i$，选择$w=i$即可，
所以式~\eqref{eq:prelim-each-sample-some-parameter}成立。但不存在一个$w$同时等于
$0$和$1$，所以式~\eqref{eq:prelim-one-parameter-all-samples}不成立。
\end{example}

例~\ref{ex:prelim-quantifier-order}说明，量词次序改变了允许依赖哪些信息。
统计学习中的“对每个允许分布，以高概率得到小误差”也包含多个量词；算法、样本量和
失败概率是否可以依赖分布，必须由量词位置决定，不能靠自然语言猜测。

\subsection{量词也规定误差界是否统一}

设$f_n:X\to\mathbb{R}$是一列函数，$f:X\to\mathbb{R}$。陈述
\[
  \forall x\in X,\quad
  \forall\varepsilon>0,\quad
  \exists N=N(x,\varepsilon),\quad
  \forall n\geq N,\quad
  |f_n(x)-f(x)|<\varepsilon
\]
允许所需的$N$依赖$x$，描述逐点收敛。若把$\exists N$移到$\forall x$之前，
\[
  \forall\varepsilon>0,\quad
  \exists N=N(\varepsilon),\quad
  \forall x\in X,\quad
  \forall n\geq N,\quad
  |f_n(x)-f(x)|<\varepsilon,
\]
同一个$N$必须同时控制所有$x$，这就是更强的一致收敛。第~\ref{chap:mathematical-analysis}
章会用具体函数序列比较两者。这里应先看到，区别不在“趋近”这个词，而在误差阈值之后
允许谁依赖谁。

\section{证明、构造与反例}
\label{sec:prelim-proofs-constructions-counterexamples}

\subsection{证明把条件连接到结论}

证明不是把若干正确公式排列在一起，而是从已给条件出发，经过每一步都合法的推理到达
结论。开始证明前，先辨认命题的逻辑形式。对$P\Rightarrow Q$，直接证明从$P$出发；
逆否证明改证$\neg Q\Rightarrow\neg P$；反证法假设$P$成立而$Q$不成立，并推出矛盾。
这些形式在逻辑上等价，但哪一种更清楚取决于条件怎样进入计算\scite{math-graham1994concrete}

\begin{theorem}[严格中点条件保证至多一个最小点]
\label{thm:prelim-strict-midpoint-unique-minimizer}
设$C\subseteq\mathbb{R}^d$满足：任取$\bm{u},\bm{v}\in C$，其中点
$(\bm{u}+\bm{v})/2$仍属于$C$。若函数$F:C\to\mathbb{R}$对任意不同的
$\bm{u},\bm{v}\in C$都满足
\begin{equation}
  F\left(\frac{\bm{u}+\bm{v}}{2}\right)
  <
  \frac{F(\bm{u})+F(\bm{v})}{2},
  \label{eq:prelim-strict-midpoint-condition}
\end{equation}
那么$F$在$C$上至多有一个最小点。
\end{theorem}

\begin{proof}
证明的关键是：若有两个不同的最小点，它们的中点会给出更小的函数值。为此反设
$\bm{u}\neq\bm{v}$都是最小点，并把共同的最小值记为$m$。由$C$对取中点封闭，
$(\bm{u}+\bm{v})/2\in C$。再由式~\eqref{eq:prelim-strict-midpoint-condition}，
\[
  F\left(\frac{\bm{u}+\bm{v}}{2}\right)
  <
  \frac{m+m}{2}
  =m.
\]
这与$m$是最小值矛盾，所以不可能存在两个不同的最小点。
\end{proof}

定理~\ref{thm:prelim-strict-midpoint-unique-minimizer}只证明“至多一个”，没有证明
最小点存在。函数可以满足排除两个最小点的条件，却根本没有最小点。第~\ref{chap:optimization}
章会把中点条件扩展为凸性，并分别建立存在条件与唯一条件。

\subsection{构造证明需要给出可检验的对象}

存在命题$\exists x\in X,\ P(x)$可以通过构造一个具体$x$并核对$P(x)$来证明。
例如，要证明对任意$a,b\in\mathbb{R}$且$a<b$，存在$c$满足$a<c<b$，可以取
$c=(a+b)/2$。由$a<b$分别加上$a$与$b$，再除以正数$2$，得到
\[
  a<\frac{a+b}{2}<b.
\]
构造不仅说明对象存在，还可能给出计算它的方法。

但存在性与可构造性不是同一个逻辑要求。某些证明通过排除“不存在”的可能来确认对象存在，
却不提供有限步骤找到该对象。反过来，一个算法给出候选对象，也仍要证明它终止并满足条件。
后文谈“最优解存在”和“算法收敛”时，会分别承担这两项责任。

\subsection{归纳法处理按整数扩展的命题}

数学归纳法适合证明由整数$n$索引的一族命题$P(n)$。它包括基例和归纳步：先证明
$P(n_0)$，再证明对每个$n\geq n_0$，$P(n)$蕴含$P(n+1)$。这两步共同覆盖所有
$n\geq n_0$，不能只展示几个数值例子。

回到式~\eqref{eq:prelim-unrolled-recurrence}，它对每个整数$t\geq1$给出同一种误差界，
正适合用归纳法验证。
\begin{proof}[式~\eqref{eq:prelim-unrolled-recurrence}的归纳证明]
当$t=1$时，式~\eqref{eq:prelim-perturbed-recurrence}正是所需结论。假设结论对$t$成立，
再使用一次递推式，得到
\begin{align*}
  e^{(t+1)}
  &\leq \rho e^{(t)}+\delta_t\\
  &\leq
  \rho^{t+1}e^{(0)}
  +
  \sum_{s=0}^{t-1}\rho^{t-s}\delta_s
  +\delta_t\\
  &=
  \rho^{t+1}e^{(0)}
  +
  \sum_{s=0}^{t}\rho^{t-s}\delta_s.
\end{align*}
这就是把$t$换成$t+1$后的式子，归纳完成。
\end{proof}

\begin{theorem}[前$n$个正整数之和]
\label{thm:prelim-arithmetic-sum}
对每个$n\in\mathbb{N}$且$n\geq1$，
\[
  \sum_{i=1}^{n}i=\frac{n(n+1)}{2}.
\]
\end{theorem}

\begin{proof}
当$n=1$时，两边都等于$1$。假设公式对某个$n\geq1$成立，则
\[
  \sum_{i=1}^{n+1}i
  =
  \left(\sum_{i=1}^{n}i\right)+(n+1)
  =
  \frac{n(n+1)}{2}+(n+1)
  =
  \frac{(n+1)(n+2)}{2}.
\]
这正是$n+1$时的公式，故结论对所有$n\geq1$成立。
\end{proof}

归纳假设只能用于已经假设的$P(n)$，不能把待证的$P(n+1)$换一种说法后当作依据。
若递推依赖前两个状态，则归纳假设和基例也要覆盖相应数量的前项。

\subsection{反例确定结论的边界}

一个全称命题声称所有允许对象都满足某种性质，所以一个合法反例就足以推翻它。反例必须
满足原命题的前提而违反结论；若它同时违反前提，只能说明例子不在讨论范围内。

例如，“目标函数有下界就一定达到最小值”是错误的。取
\[
  X=(0,1),\qquad f(x)=x.
\]
$f(x)>0$，所以$0$是下界；而对每个$x\in(0,1)$，$x/2$仍在$(0,1)$中且
$f(x/2)<f(x)$，所以没有任何$x$达到最小值。失败原因不是函数不连续，而是定义域没有
包含趋近边界时的极限点。第~\ref{chap:mathematical-analysis}章会用闭性、紧致性和连续性
给出保证极值达到的条件。

反例也能区分相近概念。一个参数化不是单射，说明参数未必可识别；它不说明预测函数一定
不能确定。一次迭代使目标下降，说明当前步骤有效；它不说明所有迭代都会收敛到全局最优。
把结论限制到证据真正覆盖的对象，是证明工作的一部分。

\section{最优值与误差表达}
\label{sec:prelim-optimal-values-errors}

\subsection{最小值与下确界}

设$X$非空且$f:X\to\mathbb{R}$。若$x^\star\in X$满足
\[
  \forall x\in X,\qquad f(x^\star)\leq f(x),
\]
则$x^\star$是一个\term{全局最小点}（global minimizer），$f(x^\star)$是最小值。
所有最小点组成
\[
  \argmin_{x\in X} f(x)
  =
  \{x\in X:f(x)=\min_{z\in X}f(z)\}.
\]
$\argmin$返回集合，而$\min$返回数值。最小点不唯一时，$\argmin$包含多个元素；
最小值未达到时，$\min$和$\argmin$都不能按上式使用。

即使没有最小值，只要$f(X)$有下界，仍可定义\term{下确界}（infimum）
\[
  f^\star=\inf_{x\in X}f(x).
\]
$f^\star$是$f(X)$的最大下界：每个$f(x)$都不小于$f^\star$，并且任何更大的数都不再
是下界。前面的$f(x)=x$、$X=(0,1)$满足$\inf_{x\in X}f(x)=0$，但不存在
$x\in X$使$f(x)=0$。

与下界对称，若$f(X)$有上界，则它的最小上界称为\term{上确界}（supremum），记为
\[
  \sup_{x\in X}f(x).
\]
若某个$x^\star\in X$达到这个数值，便得到最大值和最大点。下确界、上确界描述可逼近
的边界，最小值、最大值则额外要求边界由定义域中的元素达到。若$f$无下界，约定
$\inf_{x\in X}f(x)=-\infty$；若$f$无上界，则相应上确界为$+\infty$。

给定$\varepsilon>0$，若
\begin{equation}
  f(\widehat{x})\leq \inf_{x\in X}f(x)+\varepsilon,
  \label{eq:prelim-epsilon-optimal}
\end{equation}
则称$\widehat{x}$为$\varepsilon$-\term{近似最优解}（approximately optimal
solution）。只要下确界有限，这样的点对每个$\varepsilon>0$都存在；否则
$f^\star+\varepsilon$仍会是一个严格大于下确界的下界，与下确界的定义矛盾。
近似最优描述的是目标值差距，不自动保证$\widehat{x}$接近某个最小点。要把目标误差
换成参数距离，还需要曲率等额外条件。

式~\eqref{eq:prelim-epsilon-optimal}只适用于有限下确界。若
$\inf_{x\in X}f(x)=-\infty$，右侧不是可比较的有限误差阈值；此时应寻找一个序列
$(x_k)$使$f(x_k)\to-\infty$。例如，可以对每个正整数$k$选取满足$f(x_k)\leq-k$
的点。这样的序列说明目标值可以任意降低，却不提供一个取值为$-\infty$的“最小点”。

同一近似参数还可以用不同误差衡量。对式~\eqref{eq:prelim-regularized-least-squares}，
令$F^\star=\inf_{\bm{w}\in\mathbb{R}^d}F(\bm{w})$；若选定一个最小点
$\bm{w}^\star$，则
\[
  F(\widehat{\bm{w}})-F^\star,\qquad
  \lVert\widehat{\bm{w}}-\bm{w}^\star\rVert_2,\qquad
  \left(
    \frac{1}{n}\sum_{i=1}^{n}
    \bigl(
      \widehat{\bm{w}}^{\top}\bm{x}_i
      -(\bm{w}^\star)^{\top}\bm{x}_i
    \bigr)^2
  \right)^{1/2}
\]
分别是目标值误差、参数误差和在给定样本输入上的预测误差。最小点不唯一时，参数误差还
需要指定参照解，或改成到最小点集合的距离。三种误差衡量不同对象，彼此能否控制取决于
目标曲率、数据矩阵和参数化；一个误差较小不能无条件推出另两个也较小。

\subsection{大O、小o与比较尺度}

算法和估计误差常随规模变化。设$a_n$与$b_n$为实数序列，且从某个位置开始
$b_n>0$。若存在常数$C>0$和$n_0$，使得
\[
  \forall n\geq n_0,\qquad |a_n|\leq Cb_n,
\]
则记$a_n=O(b_n)$。这表示$a_n$的绝对值最终被$b_n$的常数倍控制，常数$C$不随$n$
改变。若
\[
  \lim_{n\to\infty}\frac{a_n}{b_n}=0,
\]
则记$a_n=o(b_n)$，表示$a_n$相对于$b_n$可以忽略。

例如$n=O(n)$、$3n+7=O(n)$且$\log n=o(n)$。但大O不是精确等式：
$n=O(n^2)$也成立，只是界较松。写算法复杂度时还要说明被计数的操作、输入规模和计算
模型；写统计误差时则要说明界是确定性的、期望意义的还是以高概率成立
\scite{math-graham1994concrete}

当序列还依赖维数$d$或其他规模时，也要说明哪些量允许进入隐含常数。例如，对每个固定
$d$都有$d/n=O(1/n)$，因为可以取$C=d$；但若要求同一个$C$同时适用于任意$d$，这个
界就不成立。因此，“关于$n$是某个速率”不等于该速率对维数一致。需要一致结论时，应把
量词写成对所讨论的全部$d$使用同一组$C$与$n_0$，或把$d$显式保留在界中。

对迭代序列，常用$t$而不是$n$表示步数。例如
\[
  F(\bm{w}^{(t)})-F^\star=O(1/t)
\]
表示存在与$t$无关的常数控制目标值误差。这个式子没有说明参数误差
$\lVert\bm{w}^{(t)}-\bm{w}^\star\rVert_2$具有同样速率，也没有说明一次迭代的成本。
比较算法时，必须固定误差对象和成本口径。

\subsection{四类问题不能由一个最优公式代替}

式~\eqref{eq:prelim-regularized-least-squares}看似只要求“求一个最小值”，实际包含四类
不同问题。

\begin{table}[htbp]
  \centering
  \small
  \begin{tabularx}{\linewidth}{@{}lX@{}}
    \toprule
    \textbf{问题} & \textbf{需要回答的内容} \\
    \midrule
    存在性 & 是否有$\bm{w}^\star$真正达到$\inf_{\bm{w}}F(\bm{w})$。 \\
    唯一性 & 达到最优值的参数是否只有一个。 \\
    可识别性 & 不同参数是否会产生相同的可观测预测或数据分布。 \\
    可计算性 & 是否存在在规定资源与精度下找到解或近似解的步骤。 \\
    \bottomrule
  \end{tabularx}
  \caption{最优化陈述中的四类不同要求。每一行需要不同的条件与论证。}
  \label{tab:prelim-four-solution-questions}
\end{table}

这四类问题彼此相关，却不能相互替代。一个有限非空候选集上的函数必定达到最小值，
但若没有给出枚举候选或计算函数值的有效方式，存在性本身不提供实际算法。例
\ref{ex:prelim-nonidentifiable-parameterization}中的参数不可唯一识别，但乘积$ab$
对应的预测函数仍可唯一确定。反过来，即使目标函数只有一个最小点，有限精度计算通常也
只能返回近似值；能否精确到达该点，或能在多大误差内逼近它，需要单独的算法保证。

回到带正则项的最小二乘。它的变量属于$\mathbb{R}^d$，目标值属于$\mathbb{R}$，
样本索引$i$遍历同一组$n$个观测，$\lambda$控制第二项的权重。若要声称某算法返回唯一
的目标最优点，需要分别证明最小点存在、最小点唯一，以及算法确实到达该点；若算法只给出
误差保证，则应改称它在规定成本下得到近似解。只有进一步声称该参数可由观测唯一确定，
或算法恢复了真实参数时，才需要额外证明数据与参数化具有所需的可识别性。后续章节将在
各自需要时分别给出这些条件。

\section{本章小结}
\label{sec:prelim-sets-functions-summary}

一个数学结论由对象、范围、条件和保证共同组成。集合与函数规定讨论什么，索引与递推规定
一族量怎样组织，量词决定选择能依赖哪些对象，证明则把条件逐步连接到结论。改变定义域、
交换量词或混淆充分条件与必要条件，都可能把原命题换成另一个命题。

最优化问题尤其需要分层判断。下确界只描述能够逼近的最佳数值，最小点还要求该数值被达到；
唯一性排除多个最小点，可识别性排除观测等价的参数，而可计算性要求给出能够在资源限制内
执行的过程。大O、小o与近似最优条件分别说明量的尺度上界、给定趋向下的相对变化和目标值
容差，不能一概理解为误差随规模缩小。第
\ref{chap:mathematical-analysis}章将沿用这些语言，研究极限、连续性、微分与积分各在什么
条件下成立。
